% !TEX program = xelatex
\ifdefined\MASTERBOOK\else
\documentclass[UTF8,zihao=-4,oneside,openany]{ctexbook}

\usepackage[a4paper,top=20mm,bottom=22mm,left=19mm,right=19mm,headheight=24pt]{geometry}
\usepackage[table]{xcolor}
\usepackage{amsmath,amssymb,mathtools,bm}
\usepackage{siunitx}
\usepackage{booktabs,tabularx,array,multicol}
\usepackage{enumitem}
\usepackage{tikz}
\usetikzlibrary{arrows.meta,positioning,calc}
\usepackage[most]{tcolorbox}
\usepackage{fancyhdr}
\usepackage{hyperref}

% ---------- Theme ----------
\definecolor{Ink}{HTML}{202326}
\definecolor{DeepBlue}{HTML}{294E69}
\definecolor{Blue}{HTML}{47687D}
\definecolor{PaleBlue}{HTML}{F3F5F6}
\definecolor{Gold}{HTML}{6E777D}
\definecolor{PaleGold}{HTML}{F6F6F5}
\definecolor{Green}{HTML}{5B6D76}
\definecolor{PaleGreen}{HTML}{F3F5F6}
\definecolor{Alert}{HTML}{666B6F}
\definecolor{PaleAlert}{HTML}{F6F6F6}
\definecolor{RuleGray}{HTML}{D3D7DA}
\definecolor{TextGray}{HTML}{61686D}

\color{Ink}
\sisetup{
  per-mode=symbol,
  inter-unit-product=\ensuremath{{}\cdot{}},
  exponent-product=\times
}

\hypersetup{
  colorlinks=true,
  linkcolor=DeepBlue,
  urlcolor=Blue,
  pdftitle={高中物理竞赛一轮讲义样章：物理量、单位与量纲},
  pdfsubject={高中物理竞赛一轮资料重排样章},
  pdfauthor={物理竞赛一轮资料编写组}
}

\setlength{\parindent}{2em}
\setlength{\parskip}{0.25em}
\setlist{nosep,leftmargin=2em}
\setcounter{secnumdepth}{3}
\setcounter{tocdepth}{2}
\numberwithin{equation}{section}
\counterwithin{table}{section}
\allowdisplaybreaks[2]
\emergencystretch=2em

\ctexset{
  chapter={
    format=\raggedright\LARGE\bfseries\color{Ink},
    titleformat=\LARGE\bfseries\color{Ink},
    beforeskip=0pt,
    afterskip=18pt
  },
  section={
    format=\Large\bfseries\color{DeepBlue},
    beforeskip=1.8ex,
    afterskip=1.0ex
  },
  subsection={
    format=\large\bfseries\color{Ink},
    beforeskip=1.5ex,
    afterskip=0.8ex
  },
  subsubsection={
    format=\normalsize\bfseries\color{Ink}
  }
}

\pagestyle{fancy}
\fancyhf{}
\fancyhead[L]{\small\color{TextGray}高中物理竞赛一轮讲义}
\fancyhead[R]{\small\color{TextGray}第一章\quad 数学与物理工具}
\fancyfoot[C]{\small\color{TextGray}\thepage}
\renewcommand{\headrulewidth}{0.4pt}
\renewcommand{\headrule}{\hbox to\headwidth{\color{RuleGray}\leaders\hrule height \headrulewidth\hfill}}

\newcommand{\dimens}[1]{\operatorname{dim}\!\left(#1\right)}
\newcommand{\dd}{\mathop{}\!\mathrm{d}}
\newcommand{\vect}[1]{\bm{#1}}
\newcommand{\tagpill}[2]{%
  \tikz[baseline=(tag.base)]\node[rounded corners=2pt,inner xsep=5pt,inner ysep=2pt,
  fill=#1!12,text=#1,font=\small\bfseries] (tag) {#2};%
}

\newtcolorbox{goalbox}{
  enhanced,breakable,
  colback=PaleBlue,colframe=Blue,
  boxrule=0.7pt,arc=2mm,
  title={本节目标},fonttitle=\bfseries,
  coltitle=DeepBlue,
  attach boxed title to top left={xshift=5mm,yshift=-2mm},
  boxed title style={colback=white,colframe=Blue,boxrule=0.7pt,arc=1.5mm}
}

\newtcolorbox[auto counter,number within=section]{examplebox}[2][]{
  enhanced,breakable,
  colback=white,colframe=RuleGray,
  boxrule=0.5pt,arc=0mm,
  title={例题~\thetcbcounter\quad #2},fonttitle=\bfseries,
  coltitle=Ink,colbacktitle=PaleBlue,
  borderline west={1.4pt}{0pt}{DeepBlue},
  #1
}

\newtcolorbox{methodbox}[1]{
  enhanced,breakable,
  colback=PaleGreen,colframe=Green,
  boxrule=0.5pt,arc=0mm,
  title={方法：#1},fonttitle=\bfseries,
  coltitle=Ink,colbacktitle=PaleGreen,
  borderline west={1.2pt}{0pt}{DeepBlue}
}

\newtcolorbox{warningbox}[1]{
  enhanced,breakable,
  colback=PaleAlert,colframe=Alert,
  boxrule=0.5pt,arc=0mm,
  title={易错提醒：#1},fonttitle=\bfseries,
  coltitle=Ink,colbacktitle=PaleAlert,
  borderline west={1.2pt}{0pt}{TextGray}
}

\newtcolorbox{extensionbox}[1]{
  enhanced,breakable,
  colback=PaleGold,colframe=Gold,
  boxrule=0.5pt,arc=0mm,
  title={拓展：#1},fonttitle=\bfseries,
  coltitle=Ink,colbacktitle=PaleGold,
  borderline west={1.2pt}{0pt}{DeepBlue}
}

\newtcolorbox{summarybox}{
  enhanced,breakable,
  colback=PaleBlue,colframe=DeepBlue,
  boxrule=0.9pt,arc=1.5mm,
  title={本节小结},fonttitle=\bfseries,
  coltitle=white,colbacktitle=DeepBlue
}

\newtcolorbox{exercisebox}[1]{
  enhanced,breakable,
  colback=white,colframe=RuleGray,
  boxrule=0.7pt,arc=1.5mm,
  title={#1},fonttitle=\bfseries,
  coltitle=DeepBlue,colbacktitle=PaleBlue
}

\begin{document}

\pagenumbering{arabic}
\setcounter{page}{1}

\chapter{数学与物理工具}
\fi
\section{物理量、单位与量纲}

物理公式不只是字母之间的代数关系。每一个字母都代表某种可以测量的物理属性，
而这种属性并不会因为单位的选取而改变。量纲正是从这一事实出发，把公式中与单位制
无关的结构提取出来。它既能用来检查推导，也能在运动规律尚不清楚时给出可靠的尺度关系。

本节从物理量的表示开始，逐步说明量纲齐次性的来源，再通过单摆、流体阻力和爆炸冲击波
等例子讨论量纲分析的能力与边界。重点不是背诵量纲表，而是学会从定义推导，并在每一步中
保留完整的量纲依据。

{\small\color{TextGray}\noindent\textbf{题源说明：}本节例题为本讲义依据量纲分析的经典模型自编，
未逐字取自特定竞赛原题；爆炸冲击波部分采用经典 Taylor--Sedov 尺度模型。\par}

\subsection{物理量不是一个孤立的数字}

物理学描述的是可比较、可测量的属性。一个物理量 $Q$ 应写成
\begin{equation}
  Q=q\,u_Q,
  \label{eq:quantity-number-unit}
\end{equation}
其中 $q$ 是在所选单位下的数值，$u_Q$ 是单位。改变单位时，物理量本身不变，数值随之改变。例如
\[
  \SI{2}{\metre}=\SI{200}{\centi\metre}.
\]
等号连接的是同一个长度，而不是数字 $2$ 与 $200$。

设同一个物理量还可以写成 $Q=q'u'_Q$。由
\[
  q\,u_Q=q'\,u'_Q
\]
可得
\begin{equation}
  q'=q\frac{u_Q}{u'_Q}.
  \label{eq:number-unit-change}
\end{equation}
这说明数值与单位的大小成反比：新单位若缩小为原来的 $1/100$，数值就扩大为原来的
$100$ 倍。物理计算中的“单位换算”并不是附加步骤，而是式 \eqref{eq:number-unit-change}
所表达的不变性。

单位给出了比较的标准，\textbf{量纲}则说明这个物理量属于哪一种物理属性。长度可以用米、厘米或天文单位表示，但它们的量纲都记作 $\mathsf{L}$。在力学中，最基本的三个量纲是
\[
  \text{长度 }\mathsf{L},\qquad
  \text{质量 }\mathsf{M},\qquad
  \text{时间 }\mathsf{T}.
\]

在更一般的国际单位制中还会用到电流 $\mathsf{I}$、热力学温度 $\Theta$、物质的量 $\mathsf{N}$ 和发光强度 $\mathsf{J}$。一轮力学问题通常只需要 $\mathsf{L}$、$\mathsf{M}$、$\mathsf{T}$。

\begin{warningbox}{数值不能脱离单位}
“这辆车的速度是 $20$”并不是完整的物理陈述。$\SI{20}{\metre\per\second}$ 与 $\SI{20}{\kilo\metre\per\hour}$ 对应完全不同的运动状态。计算过程允许暂时只处理数值，但最后结果必须恢复单位。
\end{warningbox}

\subsection{导出量纲与常见物理量}

若一个力学量 $Q$ 可以由长度、质量、时间构成，就可以写成
\begin{equation}
  \dimens{Q}=\mathsf{L}^{\alpha}\mathsf{M}^{\beta}\mathsf{T}^{\gamma},
  \label{eq:dimension-general}
\end{equation}
其中指数 $\alpha,\beta,\gamma$ 可以是整数、分数或零。量纲的推导来自物理量的定义，而不是死记硬背。例如
\[
  \dimens{v}=\frac{\mathsf{L}}{\mathsf{T}}=\mathsf{LT}^{-1},
  \qquad
  \dimens{a}=\frac{\dimens{v}}{\mathsf{T}}=\mathsf{LT}^{-2},
\]
进而
\[
  \dimens{F}=\dimens{ma}
  =\mathsf{M}\cdot\mathsf{LT}^{-2}
  =\mathsf{MLT}^{-2}.
\]
同样，由功的定义 $W=\int \vect F\cdot\dd\vect r$，积分并不改变被积量与积分变量
乘积的量纲，因此
\[
  \dimens{W}=\dimens{F}\dimens{r}
  =\mathsf{ML}^{2}\mathsf{T}^{-2}.
\]
只要沿着定义逐层展开，复杂物理量最终都会化为基本量纲的幂乘积。

\begin{table}[htbp]
\centering
\caption{一轮力学中常见物理量的量纲}
\rowcolors{2}{PaleBlue}{white}
\begin{tabularx}{\textwidth}{>{\bfseries}l >{\centering\arraybackslash}X >{\centering\arraybackslash}X l}
\toprule
物理量 & 定义或关系 & 量纲 & SI 单位 \\
\midrule
速度 $v$ & $\dd x/\dd t$ & $\mathsf{LT}^{-1}$ & \si{\metre\per\second} \\
加速度 $a$ & $\dd v/\dd t$ & $\mathsf{LT}^{-2}$ & \si{\metre\per\second\squared} \\
动量 $\vect p$ & $m\vect v$ & $\mathsf{MLT}^{-1}$ & \si{\kilogram\metre\per\second} \\
力 $F$ & $ma$ & $\mathsf{MLT}^{-2}$ & \si{\newton} \\
功、能量 $E$ & $Fs$ & $\mathsf{ML}^{2}\mathsf{T}^{-2}$ & \si{\joule} \\
功率 $P$ & $E/t$ & $\mathsf{ML}^{2}\mathsf{T}^{-3}$ & \si{\watt} \\
压强 $p$ & $F/S$ & $\mathsf{ML}^{-1}\mathsf{T}^{-2}$ & \si{\pascal} \\
角动量 $L$ & $rp$ & $\mathsf{ML}^{2}\mathsf{T}^{-1}$ & \si{\kilogram\metre\squared\per\second} \\
万有引力常量 $G$ & $Fr^2/(m_1m_2)$ & $\mathsf{M}^{-1}\mathsf{L}^{3}\mathsf{T}^{-2}$ & \si{\metre\cubed\per\kilogram\per\second\squared} \\
\bottomrule
\end{tabularx}
\end{table}

平面角是两个长度之比，因此在 SI 中属于无量纲量：
\[
  \theta=\frac{s}{r},\qquad \dimens{\theta}=1.
\]
“无量纲”不等于“没有意义”。弧度标明角度数值所采用的约定；微积分中的三角函数默认
以弧度为自变量，需要避免误读时仍应明确写出 \si{\radian}。

\subsection{量纲齐次性：公式的第一道安检}

任何有意义的物理方程都必须满足\textbf{量纲齐次性}：
\begin{enumerate}[label=\arabic*.]
  \item 等式两端量纲相同；
  \item 能够相加、相减的各项量纲相同；
  \item 指数函数、对数函数和三角函数的自变量必须无量纲。
\end{enumerate}

为什么各项必须同量纲？以只涉及长度的关系 $A=B+C$ 为例。若把长度单位从米改为厘米，
每一个长度数值都会同时乘以 $100$，方程仍保持原样。若 $C$ 不是长度量，它在换单位时就
不会按相同倍数改变，于是方程的成立将依赖于单位的选择，这显然不可能。对质量、时间等
基本量纲分别作同样的单位缩放，便得到：等式两端以及所有加减项中，每一种基本量纲的指数
都必须相同。

函数自变量必须无量纲也可以从级数看出。例如
\[
  \mathrm e^x=1+x+\frac{x^2}{2!}+\cdots,
  \qquad
  \sin x=x-\frac{x^3}{3!}+\cdots.
\]
求和各项必须同量纲，而常数 $1$ 没有量纲，所以 $x$ 只能是无量纲量。因此在振动表达式
$x=A\cos(\omega t+\varphi)$ 中，$\omega t$ 和相位 $\varphi$ 都必须无量纲。

例如匀变速运动关系
\[
  x=x_0+v_0t+\frac12 at^2
\]
中每一项的量纲都是 $\mathsf{L}$。若误写成 $x=x_0+v_0t+\frac12 at$，最后一项的量纲变成 $\mathsf{LT}^{-1}$，不可能与位移相加。

\begin{methodbox}{三步量纲检查}
\begin{enumerate}[label=\arabic*.]
  \item 把每个量替换成基本量纲 $\mathsf{L},\mathsf{M},\mathsf{T},\ldots$；
  \item 分别合并各基本量纲的指数；
  \item 比较等式两端以及每个加减项。若不一致，公式必错；若一致，公式仍需接受物理规律和边界条件的检验。
\end{enumerate}
\end{methodbox}

\begin{examplebox}{用量纲定位公式错误}
某同学把弹簧振子的机械能写成
\[
  E=\frac12 mv+\frac12 kx^2.
\]
判断错误发生在哪里。

\tcblower
弹性势能项的量纲为
\[
  \dimens{kx^2}=\dimens{\frac{F}{x}x^2}=\mathsf{ML}^{2}\mathsf{T}^{-2},
\]
确为能量量纲；但
\[
  \dimens{mv}=\mathsf{MLT}^{-1}
\]
是动量量纲。动能项应为 $\frac12 mv^2$。因此正确表达式是
\[
  E=\frac12 mv^2+\frac12 kx^2.
\]
\end{examplebox}

\begin{warningbox}{“量纲正确”只是必要条件}
公式 $s=3vt$ 与 $s=\pi vt$ 都量纲正确，但这并不能说明它们符合具体运动规律。量纲分析通常不能确定纯数系数，也不能判断正负号、矢量方向以及函数的具体形式。
\end{warningbox}

\subsection{从量纲推出尺度关系}

如果已经知道某物理量只依赖于少数几个参量，而且关系可以近似写成幂函数乘积，量纲分析就能确定各参量的指数。

一般地，设
\[
  Q=Cx_1^{a_1}x_2^{a_2}\cdots x_n^{a_n},
\]
并把每个参量写成
\[
  \dimens{x_i}=\mathsf{L}^{\ell_i}\mathsf{M}^{m_i}\mathsf{T}^{t_i},
  \qquad
  \dimens{Q}=\mathsf{L}^{\ell_Q}\mathsf{M}^{m_Q}\mathsf{T}^{t_Q}.
\]
比较三种基本量纲的指数，就得到线性方程组
\begin{equation}
  \sum_i \ell_i a_i=\ell_Q,
  \qquad
  \sum_i m_i a_i=m_Q,
  \qquad
  \sum_i t_i a_i=t_Q.
  \label{eq:dimension-linear-system}
\end{equation}
量纲分析确定幂指数，本质上就是求解这一方程组。若未知指数多于独立方程数，解中会出现
自由度；这个自由度对应某个不能由量纲决定的无量纲组合。

\begin{examplebox}{单摆周期的尺度关系}
长度为 $\ell$ 的单摆在重力加速度 $g$ 的环境中做小角度振动。假设周期 $T$ 只由 $\ell$ 和 $g$ 决定，求其尺度关系。

\tcblower
设
\[
  T=C\ell^a g^b,
\]
其中 $C$ 是无量纲常数。比较量纲：
\[
  \mathsf{T}
  =\mathsf{L}^{a}\bigl(\mathsf{LT}^{-2}\bigr)^b
  =\mathsf{L}^{a+b}\mathsf{T}^{-2b}.
\]
于是
\[
  a+b=0,\qquad -2b=1,
\]
解得 $a=\frac12$、$b=-\frac12$，所以
\begin{equation}
  T=C\sqrt{\frac{\ell}{g}}.
  \label{eq:pendulum-scaling}
\end{equation}
量纲分析确定了 $T\propto \sqrt{\ell}$ 与 $T\propto g^{-1/2}$，但不能给出 $C=2\pi$。这个系数必须由动力学方程求得。
\end{examplebox}

\begin{methodbox}{幂律关系的标准步骤}
若目标量 $Q$ 由 $x_1,x_2,\ldots,x_n$ 决定，可先假设
\[
  Q=Cx_1^{a_1}x_2^{a_2}\cdots x_n^{a_n}.
\]
把两侧写成基本量纲，对每一种基本量纲分别列指数方程。解出的 $a_i$ 给出尺度关系，$C$ 则需要通过理论、实验或边界条件确定。
\end{methodbox}

\begin{examplebox}{高速阻力的数量级}
半径为 $R$ 的球以速度 $v$ 在密度为 $\rho$ 的流体中运动。若高速情况下阻力 $F$ 主要由 $\rho$、$v$、$R$ 决定，估计 $F$ 的尺度关系。

\tcblower
设 $F=C\rho^a v^bR^c$。利用
\[
  \dimens{F}=\mathsf{MLT}^{-2},\quad
  \dimens{\rho}=\mathsf{ML}^{-3},\quad
  \dimens{v}=\mathsf{LT}^{-1},
\]
比较 $\mathsf{M}$、$\mathsf{L}$、$\mathsf{T}$ 的指数，得到
\[
  \begin{aligned}
    a&=1, & b&=2, & c&=2.
  \end{aligned}
\]
因此
\[
  F=C\rho v^2R^2.
\]
这说明高速阻力与速度平方同阶。系数 $C$ 以及公式是否适用，还取决于球的形状、表面性质和流动状态。
\end{examplebox}

\subsection{无量纲量与相似性}

两个同量纲物理量的比值是无量纲量，例如相对误差 $\Delta x/x$、折射率 $n=c/v$、相对密度 $\rho/\rho_0$。无量纲量能够把不同大小、不同单位的系统放到同一标准下比较。

表达式 $\sin(\omega t)$ 是合法的，因为 $\dimens{\omega}=\mathsf{T}^{-1}$，所以 $\omega t$ 无量纲。表达式 $\ln(t)$ 则不完整；严格写法应是
\[
  \ln\!\left(\frac{t}{t_0}\right),
\]
其中 $t_0$ 是某个参考时间。

无量纲量还揭示了量纲分析何时不能给出唯一答案。以无空气阻力的斜抛射程为例，若
$R$ 依赖 $v_0,g,\theta$，其中角度 $\theta$ 无量纲。先令
\[
  R=v_0^a g^b f(\theta),
\]
这里不能把未知的角度依赖擅自设成 $\theta$ 的某个幂，而应保留为任意无量纲函数
$f$。比较量纲：
\[
  \mathsf L=(\mathsf{LT}^{-1})^a(\mathsf{LT}^{-2})^b
  =\mathsf L^{a+b}\mathsf T^{-a-2b}.
\]
于是
\[
  a+b=1,\qquad -a-2b=0,
\]
解得 $a=2,b=-1$，故
\begin{equation}
  R=\frac{v_0^2}{g}f(\theta).
  \label{eq:projectile-dimensional-form}
\end{equation}
动力学计算会进一步给出 $f(\theta)=\sin 2\theta$，但这一步超出了量纲分析的能力。

再看球体在流体中的阻力。如果动力黏度 $\eta$ 也不可忽略，阻力应写成
\[
  F=F(\rho,v,R,\eta).
\]
四个参量只涉及三种基本量纲，因此式 \eqref{eq:dimension-linear-system} 不足以唯一确定
全部指数。可以构造无量纲的雷诺数
\[
  \mathrm{Re}=\frac{\rho vR}{\eta},
\]
于是阻力的一般量纲形式为
\begin{equation}
  F=\rho v^2R^2\,\Phi(\mathrm{Re}),
  \label{eq:drag-dimensionless-form}
\end{equation}
其中 $\Phi$ 是量纲方法无法确定的函数。高速惯性主导时，$\Phi$ 近似为常数，便回到
$F\sim\rho v^2R^2$；低速黏性主导时，$\Phi(\mathrm{Re})\sim1/\mathrm{Re}$，于是
$F\sim\eta Rv$。两种看似不同的阻力规律，因此可以统一在同一个无量纲框架中。

\begin{extensionbox}{爆炸冲击波的半径}
设一次瞬时爆炸释放能量 $E$，周围介质密度为 $\rho$。若冲击波半径 $R$ 在初始阶段只由 $E$、$\rho$ 和爆炸后的时间 $t$ 决定，令
\[
  R=CE^a\rho^bt^c.
\]
由
\[
  \dimens{E}=\mathsf{ML}^{2}\mathsf{T}^{-2},\qquad
  \dimens{\rho}=\mathsf{ML}^{-3}
\]
可列出
\[
  \begin{aligned}
    a+b&=0, & 2a-3b&=1, & -2a+c&=0.
  \end{aligned}
\]
解得 $a=\frac15$、$b=-\frac15$、$c=\frac25$，即
\[
  R=C\left(\frac{Et^2}{\rho}\right)^{1/5}.
\]
仅凭一张带时间标记的冲击波照片，就可以利用这一尺度关系反推爆炸能量的数量级。这正是量纲分析在复杂问题中的力量。
\end{extensionbox}

\subsection{单位换算与数量级检查}

单位换算的稳妥方法是把“等于 $1$ 的换算因子”逐次乘入。例如
\[
\SI{72}{\kilo\metre\per\hour}
\times\frac{\SI{1000}{\metre}}{\SI{1}{\kilo\metre}}
\times\frac{\SI{1}{\hour}}{\SI{3600}{\second}}
=\SI{20}{\metre\per\second}.
\]
每一步都能直接约去单位，不必记忆“乘还是除”。

得到数值结果后，还应做数量级和极限检查。例如地球表面附近的逃逸速度若计算成 $\SI{11}{\metre\per\second}$，即使单位正确也显然不合理；正确结果应在 $10^4\,\si{\metre\per\second}$ 的数量级。竞赛计算中，一次十秒钟的估算往往能够避免整道题因代数错误而失分。

\begin{examplebox}{从量纲关系到动力学公式}
设质量为 $M$、半径为 $R$ 的球形天体，其表面附近的第一宇宙速度 $v$ 只由
$G,M,R$ 决定。先用量纲分析求尺度关系，再由圆周运动条件确定纯数系数。

\tcblower
设
\[
  v=CG^aM^bR^c.
\]
利用 $\dimens{G}=\mathsf{M}^{-1}\mathsf{L}^3\mathsf{T}^{-2}$，有
\[
  \mathsf{LT}^{-1}
  =(\mathsf{M}^{-1}\mathsf{L}^3\mathsf{T}^{-2})^a
   \mathsf{M}^b\mathsf{L}^c
  =\mathsf{M}^{-a+b}\mathsf{L}^{3a+c}\mathsf{T}^{-2a}.
\]
逐项比较指数：
\[
  -a+b=0,\qquad 3a+c=1,\qquad -2a=-1.
\]
解得 $a=b=\frac12,c=-\frac12$，所以
\[
  v=C\sqrt{\frac{GM}{R}}.
\]
量纲分析到此只能确定尺度关系。对贴近天体表面的圆轨道，万有引力提供向心力：
\[
  \frac{GMm}{R^2}=m\frac{v^2}{R}.
\]
约去 $m$ 并整理，得到
\[
  v=\sqrt{\frac{GM}{R}},
\]
因而这个具体问题中 $C=1$。这也说明纯数系数必须由实际动力学条件补足。
\end{examplebox}

量纲检验贯穿整个物理推导：开始计算前，它帮助判断参量是否遗漏；计算过程中，它检查
每一步的加减与幂次；得到答案后，还要与数量级、极限情形和具体物理机制共同核对。
量纲不齐次的式子一定错误，而量纲齐次只说明它通过了最基本的一关。

\subsection*{分层训练}
以下各题先独立列出量纲方程，再核对参考结果。
\begin{enumerate}[label=\textbf{\arabic*.},leftmargin=2.4em]
  \item \textbf{基础} 线性阻力写成 $F=\alpha v$，平方阻力写成 $F=\beta v^2$。
  分别求 $\alpha,\beta$ 的量纲，并说明两者为什么不能使用同一个阻力系数。
  \item \textbf{基础} 已知天体表面重力加速度为 $g$、半径为 $R$。不用万有引力常量，
  由量纲关系估计贴近表面的圆轨道速度，并用 $g=\SI{9.8}{\metre\per\second\squared}$、
  $R=\SI{6.4e6}{\metre}$ 估计地球的相应速度。
  \item \textbf{方法} 密度为 $\rho$、半径为 $R$ 的小液滴发生表面振动。若周期 $T$
  只由 $\rho,R$ 和表面张力系数 $\gamma$ 决定，且
  $\dimens{\gamma}=\mathsf{MT}^{-2}$，求 $T$ 的尺度关系。
  \item \textbf{提升} 球体在流体中的阻力同时依赖 $\rho,v,R,\eta$。说明为什么量纲分析
  不能唯一确定四个参量的幂次，并写出一个含雷诺数的最简一般形式。
\end{enumerate}
{\small\color{TextGray}\noindent
参考结果：$\dimens{\alpha}=\mathsf{MT}^{-1}$，
$\dimens{\beta}=\mathsf{ML}^{-1}$；$v\sim\sqrt{gR}\approx\SI{7.9}{\kilo\metre\per\second}$；
$T\sim\sqrt{\rho R^3/\gamma}$；$F=\rho v^2R^2\Phi(\rho vR/\eta)$。\par}

\ifdefined\MASTERBOOK
  \let\sectionfinish\relax
\else
  \def\sectionfinish{\end{document}}
\fi
\sectionfinish
